% ====================================================================
\section{Introduction and aim}\label{sec:introduction}
% ====================================================================
In a static magnetic field, the two spin states of an unpaired electron have different energies (Zeeman
splitting)~\cite{weil2007,slichter1990}. Electron spin resonance (ESR) occurs when the energy of a photon, typically in
the microwave region of the electromagnetic spectrum, matches this splitting. From the resonance field one can determine
the $g$-factor of the electron. The $g$-factor is an intrinsic property of the electron in its environment and shows how
strongly orbital motion and environment change its magnetic moment. From the line width one can determine the spin-spin
relaxation time $T_2$, the time over which the spins precess in phase.

I studied three samples. DPPH (2,2-diphenyl-1-picrylhydrazyl) is a stable organic radical with one unpaired electron
per molecule. It has a narrow line and a precisely known $g$-factor and serves as the field standard~\cite{poole1996}. \ce{MnCl2} is a
concentrated paramagnet with five unpaired electrons per ion. It shows how interactions between the spins shape the
line~\cite{vanvleck1948,anderson1953}. The third sample contains \ce{Cu^2+}.

After the preparatory exercises (answered in \cref{sec:theory}), the lab manual~\cite{shiroka} sets these tasks:
\begin{enumerate}
  \item measure the magnetic field versus the control voltage, observe the hysteresis and calibrate the field
    (\cref{sec:field}),
  \item determine the field modulation amplitude and discuss its effect on the signal
    (\cref{sec:modamp}),
  \item measure the DPPH line in absorption at different modulation amplitudes and powers and estimate the line width,
    $T_2$ and $g$ from it (\cref{sec:dpph,sec:power}),
  \item measure the DPPH line in dispersion and explain the required phase shifter setting, which gives the guide
    wavelength and the cutoff frequency of the waveguide (\cref{sec:dispersion}).
\end{enumerate}
I also measured \ce{MnCl2} (\cref{sec:mncl2}) and the copper sample (\cref{sec:cu}).
